\documentclass[UTF8,11pt]{ctexart}
\usepackage{amsmath,amssymb,mathtools}
\usepackage{geometry}
\usepackage{bm}
\geometry{a4paper,margin=2.5cm}
\allowdisplaybreaks

\title{STEP2：多投入汇合与汇合后的继续传播\\
任意固定比例投入产出配方}
\author{}
\date{}

\begin{document}
\maketitle

\section{目的：汇合不终止传播}

STEP1 已经规定：

\begin{itemize}
    \item 节点自身的 production-start event 用
    \(\sigma_p,\rho_p\) 或一般的 \(D_p^{act}(t)\) 描述；
    \item 某个上游对具体下游的可见供给缺口用
    \(D_{h\to p}(t)\) 描述；
    \item production delay、transport delay 与沿途 inventory
    均在 STEP1 的 branch propagation 中处理。
\end{itemize}

STEP2 只负责回答：
\[
\boxed{
\text{多个 required inputs 已经传播到同一节点以后，
如何共同决定该节点的 production-start activity？}
}
\]

得到本地 activity 后，传播并不结束，而是重新交回 STEP1 向所有下游传播。

因此总接口为
\[
\boxed{
\text{STEP1 propagation}
\rightarrow
\text{STEP2 merge}
\rightarrow
\text{STEP1 propagation}
\rightarrow\cdots
}
\tag{1}
\]

\section{一般生产节点：任意固定比例多投入多产出}

考虑生产节点 \(P_p\)。

其必需投入集合为
\[
\mathcal I_p=\{1,\ldots,n_p\},
\]
产出集合为
\[
\mathcal O_p=\{1,\ldots,m_p\}.
\]

固定配方写为
\[
\boxed{
\sum_{h\in\mathcal I_p}a_hX_h
\longrightarrow
\sum_{r\in\mathcal O_p}b_rY_r,
}
\tag{2}
\]
其中
\[
a_h>0,
\qquad
b_r>0.
\]

这不是只针对
\[
1+1\to1,
\]
而是允许任意有限个 required inputs 与任意有限个 outputs，
并允许任意固定正投入/产出系数。

记初始化正常生产强度为
\[
\boxed{
u_p^0\equiv u_p(0).
}
\]

则第 \(h\) 种投入的初始化正常消耗速率为
\[
\boxed{
f_{p,h}^0=a_hu_p^0,
}
\tag{3}
\]
第 \(r\) 种产出的初始化正常完成速率为
\[
\boxed{
g_{p,r}^0=b_ru_p^0.
}
\tag{4}
\]

本文采用 fixed-proportion / Leontief 型 production recipe：
所有 required inputs 都不可缺少，
且暂不允许输入之间替代或动态切换配方。

\section{STEP1 交给 STEP2 的输入}

对每个
\[
h\in\mathcal I_p,
\]
设第 \(h\) 种 required input 对应一条已经传播到 \(P_p\) 输入接收平面的支路。

STEP1 返回
\[
\boxed{
D_{h\to p}(t).
}
\tag{5}
\]

其含义为：
\[
\boxed{
D_{h\to p}(t)
=
\text{截至 \(t\)，从 \(P_p\) 的视角，
第 \(h\) 种投入累计缺少的等效正常生产时间}.
}
\]

因此该投入对应的累计物料缺口为
\[
\boxed{
a_hu_p^0D_{h\to p}(t).
}
\tag{6}
\]

所有
\[
D_{h\to p}(t)
\]
必须已经位于同一个观察平面：
\[
\boxed{
\text{\(P_p\) 可以取得各 required inputs
并决定是否启动新生产的 production-start plane}.
}
\]

STEP2 不再重复计算任何已经包含在 STEP1 中的
production delay、transport delay 或 inventory coverage。

\section{单窗口输入支路与一般输入支路}

若第 \(h\) 条输入支路只有一个有效断供窗口
\[
[t_{h\to p}^{\mathrm{stop}},
t_{h\to p}^{\mathrm{rec}}),
\]
则
\[
\boxed{
D_{h\to p}(t)
=
\min\left\{
(t-t_{h\to p}^{\mathrm{stop}})_+,
\Delta_{h\to p}
\right\},
}
\tag{7}
\]
其中
\[
\boxed{
\Delta_{h\to p}
=
\left(
t_{h\to p}^{\mathrm{rec}}
-
t_{h\to p}^{\mathrm{stop}}
\right)_+.
}
\tag{8}
\]

但一般 DAG 中一条输入支路可能已经包含多个不连续缺口，
所以 STEP1 与 STEP2 之间的主接口始终是完整函数
\[
\boxed{D_{h\to p}(t),}
\]
而不是仅仅一对
\[
(t^{\mathrm{stop}},t^{\mathrm{rec}}).
\]

\section{为什么汇合不是简单的窗口并集}

如果某一种 required input 使 \(P_p\) 暂时不能生产，
其他 required inputs 也会同时停止被生产过程消耗。

这些非瓶颈投入在停产期间可能继续到货并被保存，
因而在后续时刻形成额外缓冲。

所以不能把各输入支路独立计算得到的 shortage windows
简单取并集，也不能把各自 shortage durations 直接相加。

正确对象是各输入的
\[
D_{h\to p}(t),
\]
并在 production-start plane 上比较它们的累计约束。

\section{STEP2 核心定理：required-input 上包络}

截至时刻 \(t\)，若生产节点始终以正常强度运行，
第 \(h\) 种投入应累计消耗
\[
a_hu_p^0t.
\]

由于该输入支路累计短缺
\[
a_hu_p^0D_{h\to p}(t),
\]
所以该投入截至 \(t\) 最多能够支持的累计 production activity 为
\[
\frac{
a_hu_p^0t
-
a_hu_p^0D_{h\to p}(t)
}{a_h}.
\]

于是
\[
\boxed{
\frac{
a_hu_p^0t
-
a_hu_p^0D_{h\to p}(t)
}{a_h}
=
u_p^0
\left[
t-D_{h\to p}(t)
\right].
}
\tag{9}
\]

记节点 \(P_p\) 截至 \(t\) 累计启动的 production activity 为
\[
\boxed{
U_p(t)=\int_0^t u_p(s)\,ds.
}
\tag{10}
\]

因为所有 required inputs 必须同时满足，
\[
U_p(t)
\le
u_p^0
\left[
t-D_{h\to p}(t)
\right]
\qquad
\forall h\in\mathcal I_p.
\]

在当前 greedy baseline 下，
只要物料允许就以正常强度生产，因此实际累计 activity 取最紧约束：
\[
\begin{aligned}
U_p(t)
&=
\min_{h\in\mathcal I_p}
u_p^0
\left[
t-D_{h\to p}(t)
\right]
\\
&=
u_p^0
\left[
t-
\max_{h\in\mathcal I_p}
D_{h\to p}(t)
\right].
\end{aligned}
\tag{11}
\]

定义节点自身累计 production-start 缺口
\[
\boxed{
D_p^{act}(t)
=
\max_{h\in\mathcal I_p}
D_{h\to p}(t).
}
\tag{12}
\]

于是
\[
\boxed{
U_p(t)
=
u_p^0
\left[
t-D_p^{act}(t)
\right].
}
\tag{13}
\]

在除去事件切换点后，
\[
\boxed{
u_p(t)
=
u_p^0
\left[
1-\dot D_p^{act}(t)
\right].
}
\tag{14}
\]

在完全生产/完全停产 baseline 下，
\[
\dot D_p^{act}(t)\in\{0,1\}
\]
几乎处处成立。

\section{为什么任意投入系数 \(a_h\) 不改变 \(\max\) 结构}

式 \((9)\) 表明，
每一种 required input 的数量缺口都会除以其自身配方系数 \(a_h\)，
最终统一换算成同一个量纲：
\[
\boxed{\text{等效正常 production-activity 时间}.}
\]

因此虽然
\[
a_1,a_2,\ldots,a_{n_p}
\]
可以任意不同，
归一化以后每个 required input 都给出一个
\[
t-D_{h\to p}(t)
\]
形式的 activity constraint。

所以
\[
\boxed{
\text{任意固定比例 required-input coefficients}
\quad\Longrightarrow\quad
D_p^{act}(t)=\max_hD_{h\to p}(t).
}
\tag{15}
\]

投入系数 \(a_h\) 仍然会通过正常流率
\[
f_{p,h}^0=a_hu_p^0
\]
影响物料数量与库存 coverage，
但不会改变归一化后的 \(\max\) merge 形式。

\section{任意输出系数与分支配给}

汇合以后 \(P_p\) 只有一个共同的 production-start activity
\[
u_p(t),
\qquad
D_p^{act}(t).
\]

第 \(r\) 种产出完成速率为
\[
g_{p,r}(t)=b_ru_p(t).
\]

若该产出又按固定比例
\[
\rho_{p,r\to q}
\]
分配到下游 \(P_q\)，
则该输出分支正常流率为
\[
\boxed{
f_{p,r\to q}^0
=
\rho_{p,r\to q}b_ru_p^0.
}
\tag{16}
\]

因此不同输出可以有任意固定正系数和不同固定配给比例。
这些系数缩放物料数量并进入各自分支的 inventory coverage，
但所有输出共享同一个本地 production activity
\[
D_p^{act}(t).
\]

\section{从汇合结果读取 \(P_p\) 自身的故障/恢复}

由
\[
D_p^{act}(t)
\]
读取所有满足
\[
\dot D_p^{act}(t)=1
\]
的最大连通区间：
\[
\boxed{
I_p^{act,\ell}
=
[\sigma_p^\ell,\rho_p^\ell),
\qquad
\ell=1,\ldots,L_p.
}
\tag{17}
\]

这些区间表示
\[
\boxed{
\text{\(P_p\) 自身不能启动新生产的 production-start gaps}.
}
\]

这里的
\[
\sigma_p^\ell,\rho_p^\ell
\]
仍然不是传播到任意下游后的 stop/recovery；
具体下游必须重新经过 STEP1。

若汇合结果只有一个内部窗口
\[
[\sigma_p,\rho_p),
\]
则它正好回到 STEP1 的单窗口输入形式。

\section{汇合以后如何继续冲击下游}

设 \(P_p\) 的某种产出沿一条直接分支供给下游 \(P_q\)。

STEP2 的输出不是一个终点，而是
\[
\boxed{
D_p^{act}(t),
}
\]
它立即成为下一次 STEP1 的输入。

定义输出分支的
\[
\Theta_{p\to q}
=
\tau_p^{prod}
+
\sum_{e\in\mathcal E_{p\to q}}
\tau_e^{trans},
\]
以及
\[
C_{p\to q}
=
\sum_{m\in\mathcal M_{p\to q}}
\frac{q_m(0)}{f_m^{out,0}}.
\]

则
\[
\boxed{
D_{p\to q}(t)
=
\mathcal T_{p\to q}[D_p^{act}](t)
=
\left[
D_p^{act}(t-\Theta_{p\to q})
-
C_{p\to q}
\right]_+.
}
\tag{18}
\]

因此
\[
\boxed{
\{D_{h\to p}\}_h
\overset{\mathrm{STEP2}}{\longrightarrow}
D_p^{act}
\overset{\mathrm{STEP1}}{\longrightarrow}
D_{p\to q}.
}
\tag{19}
\]

\section{单一汇合窗口下，恢复第一版的对称 downstream 公式}

若 STEP2 汇合后 \(P_p\) 只有一个自身 production-start gap
\[
[\sigma_p,\rho_p),
\]
且该 gap 足以穿透 \(P_p\to P_q\) 分支的全部库存，即
\[
\rho_p-\sigma_p>C_{p\to q},
\]
则 STEP1 立即给出
\[
\boxed{
\begin{aligned}
t_{p\to q}^{\mathrm{stop}}
&=
\sigma_p
+\Theta_{p\to q}
+C_{p\to q},
\\[1mm]
t_{p\to q}^{\mathrm{rec}}
&=
\rho_p
+\Theta_{p\to q}.
\end{aligned}
}
\tag{20}
\]

因此即使 \(P_p\) 本身是由多个 inputs 汇合得到的，
只要本地结果是一个单窗口，
对任意下游仍然完全保留
\[
\boxed{
\text{stop}=\text{local stop}+\text{delay}+\text{inventory},
\qquad
\text{recovery}=\text{local recovery}+\text{delay}.
}
\tag{21}
\]

若
\[
\rho_p-\sigma_p\le C_{p\to q},
\]
则该内部停产被这条输出分支的库存完全吸收，
\(P_q\) 不形成有效断供。

\section{多窗口汇合结果如何继续传播}

如果
\[
D_p^{act}(t)
\]
包含多个不连续增长区间，
不能把每个
\[
[\sigma_p^\ell,\rho_p^\ell)
\]
分别套用一次初始库存 \(C_{p\to q}\)。

正确做法是对整个累计函数一次性计算
\[
\boxed{
D_{p\to q}(t)
=
[D_p^{act}(t-\Theta_{p\to q})-C_{p\to q}]_+.
}
\tag{22}
\]

然后在 \(P_q\) 的接收平面读取
\[
\boxed{
[t_{p\to q}^{\mathrm{stop},j},
t_{p\to q}^{\mathrm{rec},j})
=
\text{\(\dot D_{p\to q}(t)=1\) 的第 \(j\) 个最大连通区间}.
}
\tag{23}
\]

这样同一份下游库存不会因为多个窗口而被错误地重复初始化。

\section{一个汇合节点可以对不同下游产生不同冲击}

若 \(P_p\) 同时供给
\[
P_{q_1},P_{q_2},\ldots,P_{q_s},
\]
则它们共享同一个
\[
D_p^{act}(t),
\]
但各输出分支有各自的
\[
\Theta_{p\to q_j},
\qquad
C_{p\to q_j}.
\]

因此
\[
\boxed{
D_{p\to q_j}(t)
=
\left[
D_p^{act}(t-\Theta_{p\to q_j})
-
C_{p\to q_j}
\right]_+,
\qquad
j=1,\ldots,s.
}
\tag{24}
\]

于是可能出现：
\[
P_p\text{ 自身已经停产},
\]
但某个库存充分的下游完全没有感受到断供，
而另一个库存较少的下游已经进入 shortage。

所以
\[
\boxed{
\text{local activity event 是节点共有的；
downstream-visible event 是 branch-specific 的。}
}
\tag{25}
\]

\section{两投入示例：汇合后停产、恢复、再停产}

考虑
\[
A+B\longrightarrow Y.
\]

假设 STEP1 已经把两条上游支路传播到 \(Y\) 的输入接收平面：
\[
D_{A\to Y}(t)
=
\begin{cases}
0, & t<2,\\
t-2, & 2\le t<5,\\
3, & t\ge5,
\end{cases}
\]
以及
\[
D_{B\to Y}(t)
=
\begin{cases}
0, & t<4,\\
t-4, & 4\le t<8,\\
4, & t\ge8.
\end{cases}
\]

STEP2 汇合：
\[
\boxed{
D_Y^{act}(t)
=
\max\left\{
D_{A\to Y}(t),
D_{B\to Y}(t)
\right\}.
}
\]

得到
\[
D_Y^{act}(t)
=
\begin{cases}
0, & t<2,\\
t-2, & 2\le t<5,\\
3, & 5\le t<7,\\
t-4, & 7\le t<8,\\
4, & t\ge8.
\end{cases}
\]

因此 \(Y\) 自身的 production-start gaps 为
\[
\boxed{
[2,5)
\quad\text{和}\quad
[7,8).
}
\]

若要判断这些内部 gaps 如何冲击某个下游 \(Z\)，
不能停在 STEP2，而必须继续做 STEP1：
\[
\boxed{
D_{Y\to Z}(t)
=
\left[
D_Y^{act}(t-\Theta_{Y\to Z})
-C_{Y\to Z}
\right]_+.
}
\tag{26}
\]

只有从 \(D_{Y\to Z}\) 中读出的增长区间，
才是 \(Z\) 真正观察到的
\[
t_{Y\to Z}^{\mathrm{stop},j},
\qquad
t_{Y\to Z}^{\mathrm{rec},j}.
\]

\section{与“第 \(p\) 个厂第 \(k\) 天生产状况”的对应}

STEP2 得到的
\[
D_p^{act}(t)
\]
直接决定 \(P_p\) 的生产状态：
\[
\boxed{
\frac{u_p(t)}{u_p^0}
=
1-\dot D_p^{act}(t).
}
\tag{27}
\]

若按天离散且第 \(k\) 天内部没有额外切换，
定义
\[
\boxed{
x_p[k]
=
1-
\left(
D_p^{act}(k+1)-D_p^{act}(k)
\right).
}
\tag{28}
\]

在二元 baseline 下，
\[
x_p[k]\in\{0,1\},
\]
分别表示该日正常生产或完全停产。

因此 STEP2 的 \(\max\) merge
不仅给出“汇合结果”，
还直接给出汇合节点之后原问题所关心的
\[
\boxed{\text{每个时刻/每天的 production status}.}
\]

\section{任意 DAG 的统一递归}

对任意生产节点 \(P_p\)，计算顺序固定如下。

\subsection*{输入侧：每条 required-input branch 先做 STEP1}

对每个
\[
h\in\mathcal I_p,
\]
计算
\[
\boxed{
D_{h\to p}(t)
=
\mathcal T_{h\to p}[D_h^{act}](t).
}
\tag{29}
\]

\subsection*{节点处：STEP2 做 required-input merge}

\[
\boxed{
D_p^{act}(t)
=
\max_{h\in\mathcal I_p}
D_{h\to p}(t).
}
\tag{30}
\]

并由
\[
\boxed{
u_p(t)
=
u_p^0
[1-\dot D_p^{act}(t)]
}
\tag{31}
\]
得到本地 production status。

\subsection*{输出侧：每条 downstream branch 再做 STEP1}

对每个直接下游 \(P_q\)，
\[
\boxed{
D_{p\to q}(t)
=
\mathcal T_{p\to q}[D_p^{act}](t).
}
\tag{32}
\]

因此整个 DAG 只反复做两类操作：
\[
\boxed{
\text{沿 branch：delay shift + inventory absorption},
}
\]
以及
\[
\boxed{
\text{到 production node：required-input \(\max\) merge}.
}
\]

可以写成
\[
\boxed{
\mathcal T
\rightarrow
\max
\rightarrow
\mathcal T
\rightarrow
\max
\rightarrow\cdots.
}
\tag{33}
\]

\section{与单输入节点的一致性}

若
\[
|\mathcal I_p|=1,
\]
则
\[
\boxed{
D_p^{act}(t)
=
D_{h\to p}(t).
}
\tag{34}
\]

因此 STEP2 自动退化为恒等映射。
这时整个网络沿直链只剩 STEP1 的算子复合，
也就退化回第一版的逐段累加公式。

\section{库存归属与禁止重复计算}

必须遵守
\[
\boxed{
\text{一份物理库存只能属于一个明确的 propagation segment}.
}
\]

如果某库存已经计入某条输入支路的
\[
C_{h\to p},
\]
则 STEP2 不能再次把它作为新的 buffer 扣除。

如果 \(P_p\) 自身具有独立 input-side local storage，
应在拓扑中明确表示为库存节点 \(M\)，
并明确归属于哪一条 STEP1 propagation segment。

同理，
finished-goods storage 若要吸收 \(P_p\) 的输出缺口，
也应作为输出侧 STEP1 中的库存节点处理，
不能隐藏进 production delay 或 STEP2 merge。

\section{当前 \(\max\) 闭式规则的适用边界}

式
\[
D_p^{act}(t)=\max_hD_{h\to p}(t)
\]
基于以下 baseline：
\begin{itemize}
    \item 所有投入为固定比例的 required inputs；
    \item 输入之间不能直接替代；
    \item 生产配方 \(a_h,b_r\) 固定；
    \item 正常情况下只要物料允许，节点以 \(u_p^0\) 运行；
    \item 一个投入导致共同停产时，其他到货投入可以被保存；
    \item 暂不考虑保存容量上限、损耗、报废与保质期；
    \item 不主动补库、不超额生产、不重新路由；
    \item 当前主要考虑完全停产---完全恢复。
\end{itemize}

因此这里的“任意投入产出比”准确含义是
\[
\boxed{
\text{任意固定正系数的 multi-input / multi-output Leontief recipe},
}
\]
而不是任意可替代、可切换或非线性的 production technology。

若非瓶颈投入在共同停产期间无法继续保存，
例如 local storage capacity 绑定，
则简单的
\[
D_p^{act}(t)=\max_hD_{h\to p}(t)
\]
可能失效，需要回到显式库存守恒方程。

若允许 partial shortage，
可以推广为
\[
0\le\dot D(t)\le1,
\]
此时
\[
u_p(t)
=
u_p^0
[1-\dot D_p^{act}(t)]
\]
自然给出部分生产。

\section{STEP2 最终接口}

STEP2 可以压缩为
\[
\boxed{
\underbrace{
\{D_{h\to p}(t)\}_{h\in\mathcal I_p}
}_{\text{STEP1 已传播到 \(P_p\) 的 required-input shortages}}
\quad
\xrightarrow{\ \max\ }
\quad
\underbrace{
D_p^{act}(t)
}_{\text{\(P_p\) 本地 production-start activity gap}}
}
\tag{35}
\]

其中
\[
\boxed{
D_p^{act}(t)
=
\max_{h\in\mathcal I_p}
D_{h\to p}(t).
}
\]

随后对每一个直接下游 \(P_q\)：
\[
\boxed{
D_p^{act}(t)
\quad
\xrightarrow{\ \mathrm{STEP1}\ }
\quad
D_{p\to q}(t)
=
[D_p^{act}(t-\Theta_{p\to q})-C_{p\to q}]_+.
}
\tag{36}
\]

所以最终闭环是
\[
\boxed{
\text{upstream local activity}
\rightarrow
\text{STEP1 propagation}
\rightarrow
\text{required-input shortages}
\rightarrow
\text{STEP2 max merge}
\rightarrow
\text{local production activity}
\rightarrow
\text{STEP1 propagation}.
}
\]

单窗口时可随时退回第一版的对称公式；
多窗口时保留完整 \(D(t)\) 继续递归。

\end{document}
